\section{Valuative pushforward of discrepancy functions}
\label{b:discrepancy-push}

The residual system from Section~\ref{a:essential} will be studied through
successively smaller function fields. We need discrepancy data that pass
through these contractions with an exact rule for pullbacks and with
controlled denominators. The construction below uses the least normalized
discrepancy among divisors lying over a given base divisor. We first prove
attainment and compatibility, then construct its nef moduli presentation.

We continue over \(\C\), and retain the chosen rational volume forms.
A prime divisor is identified with its normalized integral valuation.
The coefficient of a rational b-divisor \(\mathbf D\) at a valuation
\(P\) is denoted by \(P(\mathbf D)\). All b-divisors are evaluated on
models on which the displayed Cartier data descend.

\begin{definition}\label{b:discrepancy-presentation}
A discrepancy presentation on a normal projective model \(Y\) is
\[
                 a=a_{K,Y}+\mathbf M-\overline D.
\]
Here \(a_{K,Y}\) is the form discrepancy, \(\mathbf M\) is b-nef and
b-\(\Q\)-Cartier, and \(D\) is a rational Cartier divisor on \(Y\).
The boundary trace \(B=\sum_P(1-a(P))P\), summed over strict primes
on \(Y\), satisfies
\[
                       K_Y+B+M_Y=D
\]
as divisors. The presentation is \emph{Calabi--Yau} if \(D=0\).
Sub-generalized lc and sub-generalized klt mean, respectively,
\(a(P)\ge0\) and \(a(P)>0\) for every relevant prime divisor on every
higher model. Effectivity means effectivity of the boundary trace
on the specified model, not on every higher model.
\end{definition}

The effective unscaled case is the generalized-pair framework of
Birkar--Zhang \cite[Definitions~1.4 and~4.1]{BirkarZhang}.
Positive rational scalings are also allowed. For
\(\mu a_K+\mathbf M-\overline D\), trace effectivity means that strict
prime discrepancies are at most \(\mu\). Bars denote Cartier closure
and may be suppressed when there is no ambiguity.

\begin{definition}\label{b:valuative-push}
For a projective contraction \(\pi:Y\to V\), put
\begin{equation}\label{eq:B1}
 (\pi_\#a)(P)=
       \min_{\substack{F|_{\C(V)}=mP\\m>0}}\frac{a(F)}m.
\end{equation}
We use this operation for data sub-generalized lc over the generic
point of \(V\), with total divisor pulled back from \(V\), and for
positive scalings of such data. A prime valuation upstairs is
\emph{horizontal} when its restriction to \(\C(V)\) is trivial.
\end{definition}

For an unscaled discrepancy presentation, \((\pi_\#a)(P)\) is one
minus the discriminant coefficient at \(P\) in the canonical bundle
formula \cite[Section~2]{FilipazziGeneralized}. Keeping it as a
function on all prime valuations makes composition and base pullbacks
explicit.

\begin{lemma}\label{b:push-threshold}
The minimum in \eqref{eq:B1} is finite and attained. Pullback terms
from the base translate it exactly, and successive pushes compose.
\end{lemma}

\begin{proof}
Place \(P\) on a smooth projective base model. Resolve the map upstairs,
the boundary and nef data, and the pullback of \(P\), so that the nef
part descends and the relevant divisors have simple normal crossings.
Over \(\eta_P\), the largest real number that can be subtracted times
the pullback of \(P\) while preserving sub-lc discrepancies is the
least component-discrepancy-to-multiplicity ratio among components
of positive multiplicity. This number may be negative. Horizontal
sub-lc takes care of components of zero multiplicity. The snc
discrepancy formula gives the same inequalities for every higher
valuation, proving attainment and \eqref{eq:B1}.

A nontrivial restriction of a prime-divisor valuation to a subfunction
field is a positive integral multiple of a prime-divisor valuation:
the relative Abhyankar inequalities force equality for the restriction,
and its nonzero value group is a subgroup of \(\Z\).
Apply this to the intermediate fields and choose attained minimizing
lifts successively to obtain composition.
Finally \(F(\pi^*D)=mP(D)\), proving translation.
\end{proof}

\subsection{Preparatory birational operations}

\begin{lemma}\label{b:ordinary-approximation}
Effective generalized klt data with big b-nef part admit an ordinary
effective klt representative in the same rational adjoint class, whose
discrepancies do not exceed the generalized ones. Generalized
\(\delta\)-lc data, \(0<\delta\le1\), admit such a representative which
is \(\delta/2\)-lc. A sufficiently small ample class on the low model
may be reserved before taking the representative. In particular,
effective relatively Calabi--Yau generalized klt data with big b-nef
part make the underlying variety of Fano type over the base.
\end{lemma}

\begin{proof}
On a projective descent write \(M\sim_\Q D_0+E_0\), with \(D_0\)
ample and \(E_0\ge0\). For small rational \(s>0\), use \(sE_0\) as
fixed part and the ample class \((1-s)M+sD_0\) as moving part.
A still smaller pullback ample reserve can be subtracted from this
moving part. Resolve fixed supports and take a sufficiently divisible
general free representative of the moving part, divided by its
multiplicity. Fixed coefficients change arbitrarily little and new
coefficients are arbitrarily small. The snc discrepancy formula
therefore retains klt, or \(\delta/2\)-lc, on every higher valuation.
Subtracting the effective Cartier b-divisor representing the nef part
only lowers the generalized discrepancies. For open bases use a
projective compactification and common higher projective descents,
testing these inequalities over the required open. The ample reserve
proves the Fano-type assertion.
\end{proof}

On a \(\Q\)-factorial model the negativity lemma
\cite[Lemma~4.16]{FujinoFundamental} gives
\[
                         \overline{M_Y}-\mathbf M\ge0.
\]
Removing effective boundary or forgetting b-nef data consequently
does not decrease discrepancies. Adjoint MMPs preserve lower
discrepancy bounds by pullback negativity on common resolutions.

We use the following established model results with their hypotheses.
An effective klt rational pair with big boundary and pseudo-effective
adjoint, projective over a quasi-projective base, has a log terminal
model; the model is nonextracting and \(\Q\)-factorial, and is good
because the boundary is big
\cite[Theorem 1.2 and Lemma 3.9.3]{BCHM}.
The nonextraction convention is explicit in
\cite[Definition~3.6.6]{BCHM}; all these locators refer to the same
arXiv version.
Relative bigness is tested on the generic fiber. A rational-function
witness can be extended globally after adding a sufficiently vanishing
effective Cartier divisor from the base to remove its finitely many
residual vertical poles.

A projective \(\Q\)-factorial variety of Fano type is a Mori dream
space \cite[Corollary 1.3.1]{BCHM}; in particular, its effective cone
is rational polyhedral and closed, rational classes in it are
rational-linearly effective, and their divisor MMPs end at nef
semiample models \cite[Definition 1.10 and Proposition 1.11]{HuKeel}.
This is a statement for each individual variety, with no uniform
chamber or generator assertion. Nonextracting \(\Q\)-factorial
outputs from a Fano-type model remain of Fano type: push an effective
klt rationally trivial adjoint with big boundary. Rational triviality
makes this crepant, and persistence of sections preserves bigness.
Perturbing the big boundary gives a log Fano boundary again.

\begin{lemma}\label{b:controlled-extraction}
Let \((Y,B)\) be an effective klt log Fano pair. Any finite collection
of exceptional prime divisors with discrepancies in \((0,1]\) can be
extracted on a projective \(\Q\)-factorial birational model having no
other exceptional prime divisors. This model is of Fano type.
An empty collection gives a small \(\Q\)-factorialization.
\end{lemma}

\begin{proof}
On a log resolution displaying the collection, retain their crepant
coefficients, which lie in \([0,1)\). Raise every other exceptional
coefficient to a strictly larger nonnegative rational number below
one. The resulting effective pair is klt, its boundary is big over
the birational base, and its adjoint is relatively the effective
exceptional divisor \(R\) supported exactly on the raised components.
Normality gives \(f_*\OO(mR)=\OO_Y\) for divisible \(m\).
The good log terminal model preserves these sections by negativity
and makes the transform of \(R\) relatively semiample. A free multiple
has only the relative generator \(1\), so its effective transform
vanishes. The pullback difference is zero at every requested divisor;
strict adjoint negativity for contracted divisors therefore prevents
their contraction. The MMP extracts no divisor. The original pair's
crepant trace on this extraction is effective and has nef big
anti-adjoint, so a small perturbation gives Fano type.
\end{proof}

\subsection{A uniform interval of primary minimization}

For a small perturbation of an lc discrepancy function, we need its
computing divisors to continue computing the unperturbed function. The
next lemma obtains one interval valid for every tested base prime from a
bounded index and generalized threshold ACC. This is the uniformity that
will later allow us to pass to a nef endpoint on a fixed model.

\begin{lemma}\label{b:primary-minimization}
Fix a positive integer \(r_{\mathrm{ind}}\). Let \(A,C\) have Calabi--Yau presentations on \(Y\), effective in
horizontal trace over \(V\). Assume that \(A\) is sub-generalized lc
over the generic point, that \(C\) is generalized klt there with big
b-nef part, and that \(r_{\mathrm{ind}}\mathbf M_A\) is integral b-Cartier.
There is \(t_0>0\), depending only on the dimension and \(r_{\mathrm{ind}}\), such
that for every rational \(0<t<t_0\), every prime \(T\) computing
\(\pi_\#((1-t)A+tC)(P)\) also computes \(\pi_\#A(P)\).
\end{lemma}

\begin{proof}
Fix a prime \(P\) and a lift \(T\) computing the mixture. We will
construct a model on which \(T\) is the only strict prime above \(\eta_P\), and
on which both individually shifted boundaries are effective. On that
model, generalized threshold ACC will force the shifted \(A\)-datum
to be lc.

\medskip\noindent
\emph{Shifting the two data at the computing prime.}
Work over an open neighborhood \(U\) of \(\eta_P\) on a smooth base
model and put
\[
 T|_{\C(V)}=m_TP,\qquad
 \alpha_A=A(T)/m_T,\qquad \alpha_C=C(T)/m_T.
\]
Subtract \(\alpha_A\pi^*P,\alpha_C\pi^*P\) from the respective
discrepancy functions. Their mixture is sub-generalized lc over \(U\)
after shrinking: its threshold at \(P\) is zero, and horizontal sub-lc
is already known. Both shifted totals are pulled back from \(U\).

Add \(\kappa\pi^*P\) to the shifted mixture discrepancy for a small
positive rational \(\kappa\). It becomes sub-generalized klt there,
with coefficient \(1-\kappa m_T\ge0\) at \(T\). On a high smooth
model displaying \(T\), resolve the boundary, exceptional, fiber and
nef data. Raise the coefficients of every other component over
\(\eta_P\), and every horizontally exceptional divisor over the
original generic fiber, to strictly larger nonnegative rationals
below one. Remove unwanted vertical supports away from \(\eta_P\)
by shrinking. Horizontal strict primes over the original generic
fiber already have effective trace. Thus we obtain effective
generalized klt data whose adjoint is relatively an effective divisor
\(R_{\mathrm{ex}}\), supported exactly on these raises and omitting
\(T\).

\medskip\noindent
\emph{A model whose only strict prime above \(\eta_P\) is \(T\).}
By Lemma~\ref{b:ordinary-approximation}, replace the big nef part by
ordinary effective data while keeping klt and big boundary.
Take its good \(\Q\)-factorial log terminal model over \(U\).
The reason every raised divisor disappears is that
\(R_{\mathrm{ex}}\) has no relative sections except those from
\(\OO_U\). We verify this before using semiampleness on the good
model.

For every divisible \(m\), a relative section of
\(\OO(mR_{\mathrm{ex}})\) has no nonexceptional horizontal poles
on the original normal proper generic fiber. It is therefore a
constant on that fiber, hence an element of \(\C(V)\).
A base pole at \(P\) is impossible because its pullback has a pole
along \(T\), absent from \(R_{\mathrm{ex}}\). At every other retained
base prime, a component missing from \(R_{\mathrm{ex}}\) likewise
prevents a pole. Such a component exists above each strict prime of
a normal contraction base, by pulling back a local equation near its
generic point; other vertical supports were removed by the shrink.
Consequently these relative section sheaves are precisely \(\OO_U\).

The ordinary approximation may change the adjoint by a rational
principal term and a base pullback. Clear denominators and identify
the corresponding section spaces by multiplication by that principal
function and local trivialization of the base term. Under this
identification the algebra just computed is unchanged. Negativity
preserves it on the good model. The effective transform of
\(R_{\mathrm{ex}}\) is relatively semiample; a divisible free multiple
is generated by the section \(1\), whose zero divisor is that multiple
of \(R_{\mathrm{ex}}\). Hence its support is empty.
The output must still have a divisor over \(\eta_P\), and the MMP
extracts none. Thus \(T\) survives and is the unique such divisor.

The pre-raise, \(\kappa\)-adjusted mixture is crepant on this output,
since its adjoint is pulled back from \(U\), and has effective trace.
Its big b-nef part gives relative Fano type by
Lemma~\ref{b:ordinary-approximation}. In particular the bare output
is klt. The surviving horizontal primes come from the original
generic fiber; both individual shifted boundaries are therefore
effective after shrinking. Their coefficient at \(T\) is one.

\medskip\noindent
\emph{The bounded-coefficient threshold test.}
We now return to the shifted mixture before the \(\kappa\)-adjustment;
it is lc on this same model. Remove \(t\) times the shifted
\(C\)-boundary and its corresponding nef datum. Effectivity and
negativity show that discrepancies only increase. Thus the
generalized threshold on the bare variety of the
shifted \(A\)-boundary \(B'_A\) together with \(\mathbf M_A\) lies in
\([1-t,1]\); the upper bound comes from \(T\).
The horizontal coefficients of \(B'_A\) lie in the finite
\(1/r_{\mathrm{ind}}\)-grid in \([0,1]\), because \(r_{\mathrm{ind}}A(F)\) is integral for Calabi--Yau
data with \(r_{\mathrm{ind}}\mathbf M_A\) Cartier. The only remaining nonzero vertical
coefficient is one. Use \((1/r_{\mathrm{ind}})(r_{\mathrm{ind}}\mathbf M_A)\) as a single nef
Cartier summand. These are exactly the effective test and nef
Cartier coefficient hypotheses of generalized threshold ACC
\cite[Theorem 1.5]{BirkarZhang}. The test scales the boundary and the
nef datum together; at parameter one it is precisely the shifted
\(A\)-presentation. More explicitly, the initial boundary and nef
datum in the threshold theorem are both zero; the effective test
boundary is \(B'_A\), and the nef test datum is
\((1/r_{\mathrm{ind}})(r_{\mathrm{ind}}\mathbf M_A)\).
The latter is required to be nef, with the stated Cartier coefficient;
its trace need not be effective. A Cartier divisor descending on the low model
does not affect the nef-data discrepancy
\(\mathbf M_A-\overline{M_{A,Y}}\), so any auxiliary effective
representative obtained by twisting on the affine base does not
alter the test.

If thresholds below one occurred for arbitrarily small \(t\), a
subsequence would increase strictly to one, contradicting ACC.
Thus the threshold is one for all sufficiently small \(t\), with a
bound depending only on dimension and \(r_{\mathrm{ind}}\). Closedness supplies the
inequalities at the endpoint. The shifted \(A\)-data are therefore
lc over \(\eta_P\), proving that \(T\) computes \(\pi_\#A(P)\).
\end{proof}

\begin{corollary}\label{b:principal-denominators}
In Lemma~\ref{b:primary-minimization}, with the fixed positive integer
\(r_{\mathrm{ind}}\) as there, assume in addition that
\(r_{\mathrm{ind}}\mathbf M_A\) is principal. Then every value \((\pi_\#A)(P)\)
has bounded denominator, in terms only of the dimension and \(r_{\mathrm{ind}}\).
\end{corollary}

\begin{proof}
Use the model constructed in the proof of
Lemma~\ref{b:primary-minimization}: its unique strict prime above \(\eta_P\) is
\(T\), and the shifted \(A\)-boundary is effective and lc. The
shifted \(A\)-data now have ordinary discrepancies, and
\[
                 K+B'_A=\alpha_A\pi^*P-M_{A,Y}.
\]
The pair is lc, its anti-adjoint is nef over \(U\), its coefficient
set is finite rational (hence contained in a fixed hyperstandard
set), and its underlying variety is of Fano type over \(U\).
Apply the bounded relative monotone lc complement theorem near a
closed point of \(P\cap U\)
\cite[Theorem 1.8]{BirkarComplements}.
The effective complement addition restricts to a rationally trivial
effective divisor on the proper generic fiber, so has no horizontal
component. It cannot add to \(T\), whose coefficient is already one.
Thus the boundary is unchanged near \(\eta_P\). Higher crepant
traces are determined by the same adjoint and introduce no new
independent addition.

For clarity, if the local complement identity is
\(n'(K+B'_A)=\Div(h)\) and
\(r_{\mathrm{ind}}M_{A,Y}=\Div(g)\), take
\(N=n'r_{\mathrm{ind}}\) and \(\phi=h^{r_{\mathrm{ind}}}g^{n'}\).
Thus a bounded integer \(N\) gives
\[
                   N\alpha_A\pi^*P=\Div(\phi)
\]
over a neighborhood still containing \(\eta_P\).
The principal adjustment from \(M_A\) is absorbed in \(\phi\).
This function has no horizontal poles. Generic-fiber properness and
the contraction property imply \(\phi\in\C(V)\), even though its
initial expression was obtained on an open neighborhood.
Evaluating at \(T\) gives
\[
           N\alpha_A m_T=m_T\ord_P(\phi),
             \qquad N\alpha_A=\ord_P(\phi)\in\Z.
\]
The fiber multiplicity cancels; neither a norm nor an extension-degree
factor is introduced.
\end{proof}

\subsection{Nef push presentations}

The preceding denominator argument applies to principal moduli data.
We now construct the moduli b-divisor for a general semiample or nef and
big presentation. The proof treats descent from generic coefficients,
tests on newly appearing primes, and nefness separately, before taking
the endpoint of a small perturbation. For the unscaled presentation, the
moduli conclusion is a case of the generalized canonical bundle formula
of Filipazzi \cite[Theorem~1.4]{FilipazziGeneralized}. We retain an
explicit proof in the chosen volume-form conventions, together with the
perturbation argument used here.

\begin{proposition}\label{b:push-presentation}
Let \(a=a_{K,Y}+\mathbf M\) be a Calabi--Yau presentation, sub-generalized
klt over the generic point of \(V\) and effective in horizontal trace.
If \(\mathbf M\) is b-semiample or b-nef and big, then
\[
                        \pi_\#a=a_{K,V}+\mathbf N
\]
for a b-nef b-\(\Q\)-Cartier divisor \(\mathbf N\).
Positive rational scaling and translation by base Cartier pullbacks
give the corresponding presentations with their scaled canonical
coefficient and prescribed total divisor.
\end{proposition}

\begin{proof}
We first assume b-semiampleness. Let \(k=\C\), and enlarge constants
to the algebraic closure \(k'\) of a purely transcendental extension
of finite transcendence degree. Use enough independent coefficients
to choose a member of a free Cartier system \(|m\mathbf M|\),
\(m>1\), generic over \(k\), on a high descent. Divide it by \(m\)
and denote its Cartier closure by \(\Delta\).
The data
\[
                       a^{k'}-\Delta
        =a_{K,Y}^{k'}+\mathbf M^{k'}-\Delta
\]
are ordinary pair discrepancies, because the remaining moduli term
is rationally principal. Their ordinary boundary adds the trace of
\(\Delta\) to the generalized boundary. Bertini on a log resolution,
with moving coefficient \(1/m<1\), makes them sub-klt generically.
Their horizontal trace on the original low model remains effective.

The rank-one condition in the ordinary canonical bundle formula
holds as follows. On a resolution over the generic point, the
rounded discrepancy sheaf contains constants by sub-klt singularities.
Any additional poles allowed by it are exceptional over the original
normal generic fiber, because the low horizontal boundary is
effective. Such functions are regular on that proper low fiber and
therefore constant. The contraction condition gives rank one.
The rational principal adjustment makes the total adjoint pulled back
from zero. Ambro's ordinary klt-trivial-fibration theorem therefore
gives a b-nef b-\(\Q\)-Cartier divisor
\[
 \mathbf N^\Delta
    =(\pi^{k'})_\#(a^{k'}-\Delta)-a_{K,V}^{k'}.
\]
Here we use \cite[Theorem 0.2]{AmbroBoundary}; see also
\cite{AmbroModuli}.
Vertical negative boundary components are allowed in this theorem.
Alternatively, compensate them with a pulled-back base divisor,
which translates thresholds and total divisor compatibly.

For each prime \(P\) over \(k\), a resolution computing its threshold
is defined over \(k\). Extension does not change its finite snc
calculation, and the generic free divisor meets that calculation
transversely, without changing the threshold. Thus, writing
\(\mathbf N=\pi_\#a-a_{K,V}\) initially as a coefficient function,
\begin{equation}\label{b:old-prime-equality}
                    \mathbf N^\Delta(P_{k'})=\mathbf N(P).
\end{equation}
This holds for every old prime simultaneously in the requisite sense:
each transversality condition is a nonempty open condition defined
over \(k\), and the chosen coefficient tuple lies over the generic
point of its parameter space. No countable intersection of open
subsets of \(k\)-points is being asserted.

\medskip\noindent
\emph{Descent of the Cartier presentation.}
The restriction to old prime extensions of any rational b-Cartier
divisor over \(k'\) is rational b-Cartier over \(k\).
On a projective descent, taking differences and multiples reduces
this to a globally generated Cartier divisor \(J\). For finitely many
rational sections \(f_j\),
\[
                          Q(J)=-\min_j Q(f_j).
\]
Put these sections over a common denominator in
\(k(V)\otimes_k k'\), and express every numerator and the denominator
as a finite sum of coefficients in \(k(V)\) times constants linearly
independent over \(k\). At \(P_{k'}\), the value of such a sum is
exactly the least coefficient value. Indeed those constants remain
independent over \(k(P)\) in
\(\Frac(k(P)\otimes_k k')\); geometric integrality over algebraically
closed \(k\) makes this tensor product a domain.
Hence the restricted function is a difference of minima over finite
lists in \(k(V)\). Resolving the corresponding fractional ideals
makes it Cartier. Apply this to \(\mathbf N^\Delta\) and
\eqref{b:old-prime-equality} to obtain a b-\(\Q\)-Cartier presentation
of \(\mathbf N\) over \(k\).

\medskip\noindent
\emph{Testing new primes and nefness.}
The identity \(\pi_\#a=a_{K,V}+\mathbf N\) remains valid after
extension at every prime over \(k'\). To prove this, fix one such
prime on a smooth model dominating a descent of \(\mathbf N\).
Choose a smooth model upstairs, with descended \(\mathbf M\),
on which its threshold is computed by snc data including the
pullback of that prime. All models, morphisms, prime components,
rational forms and Cartier data in this single test are defined over
a finitely generated extension of \(k\). Spread them over an
integral parameter variety over \(k\). After shrinking, the models
remain projective birational resolutions fiberwise; the required
components remain geometrically integral; smoothness and relative
snc hold; and the same components dominate the tested prime.
Properness and constructibility permit these requirements.
The finitely many divisor identities for forms, functions and
pullbacks also persist with unchanged multiplicities: keep their
supports on the smooth models and remove the parameter locus where
a stated divisor identity or dominance fails.

Specialize to a sufficiently general closed \(k\)-point. Both the
proposed b-Cartier coefficient and the threshold are computed by
the same respective finite numerical lists as before specialization.
They agree there by the already proved identity over \(k\).
They therefore agree for the chosen prime over \(k'\).
This argument uses one finite spread per test prime, not a common
resolution of all primes.

Since \(\Delta\ge0\), threshold monotonicity gives the b-effective
difference
\[
                         \mathbf N^{k'}-\mathbf N^\Delta\ge0.
\]
Its value at every \(P_{k'}\) is zero by
\eqref{b:old-prime-equality}. Let \(C\) be a curve on a smooth
projective descent of \(\mathbf N\) over \(k\), and take a common
higher descent over \(k'\) where \(\mathbf N^\Delta\) is nef.
Choose an old prime valuation centered at the generic point of \(C\).
Its extended center \(Z\) on the common model dominates \(C_{k'}\).
After clearing denominators, if \(Z\) lay in the effective Cartier
difference, a local equation would have strictly positive valuation,
contrary to its zero value. Thus \(Z\) is not contained in that
support. General complete intersections in \(Z\) have a component
dominating \(C_{k'}\) and not contained in the support: their
restriction to the generic fiber of \(Z\to C_{k'}\) has positive
top ample intersection, and a general choice avoids containment in
a fixed proper closed subset. On this curve the effective difference
has nonnegative degree. The nef term and projection formula show
\(\mathbf N\cdot C\ge0\). If the base itself is a curve, use the
whole curve on its common smooth model; the avoidance issue is
automatic. A zero-dimensional base is immediate.
This proves b-nefness in the semiample case.

\medskip\noindent
\emph{The nef and big endpoint.}
On a projective descent write \(M\sim_\Q D_0+E_0\), with \(D_0\)
ample and \(E_0\ge0\). For sufficiently small rational \(s>0\),
\[
                           a_s=a-s\overline E_0
\]
retains horizontal klt and effectivity, and its moduli class
\(M-sE_0\sim_\Q(1-s)M+sD_0\) is ample, hence semiample.
Fix one such \(s_0>0\). Apply Lemma~\ref{b:primary-minimization}
to \(a,a_{s_0}\), choosing a positive integer \(r_{\mathrm{ind}}\) for
which \(r_{\mathrm{ind}}\mathbf M\) is integral b-Cartier. This index may
depend on the individual presentation; the interval below is uniform over
its test primes.
Since
\[
 a_s=(1-s/s_0)a+(s/s_0)a_{s_0},
\]
the lemma gives one interval \(0<s<s_1\), independent of the
prime \(P\), on which every computing lift for
\(\pi_\#a_s(P)\) also computes \(\pi_\#a(P)\).

For a fixed \(P\), a computing lift at any one positive parameter
exists by Lemma~\ref{b:push-threshold}. Comparing it with every
other primary minimizer shows that it maximizes
\(E_0(F)/m\) among them. This maximum is attained and is independent
of \(s\). Thus the threshold is exactly
\[
                    (\pi_\#a_s)(P)
                         =(\pi_\#a)(P)-s e_P
\]
on this interval, including its stated value at \(s=0\).
The identity holds for every prime. If \(0<s'<s''<s_1\) are
rational, the endpoint moduli function is therefore
\[
 \mathbf N_0=
       \frac{s''\mathbf N_{s'}-s'\mathbf N_{s''}}{s''-s'}.
\]
The two terms on the right descend on a common model; hence so does
\(\mathbf N_0\). All other small positive-parameter terms descend
there by the same affine formula and are nef there, since nefness
on a higher descent descends by lifting curves and projection
formula. Taking the limit on this fixed model proves that
\(\mathbf N_0\) is nef. This proves the endpoint assertion.
Scaling and translation now follow from
Lemma~\ref{b:push-threshold}.
\end{proof}
